\documentclass[11pt]{article}

\usepackage[a4paper,margin=1in]{geometry}
\usepackage[T1]{fontenc}
\usepackage[utf8]{inputenc}
\usepackage{lmodern}
\usepackage{microtype}
\usepackage{amsmath,amssymb,amsthm,mathtools}
\usepackage{booktabs}
\usepackage{xcolor}
\usepackage[colorlinks=true,linkcolor=blue!45!black,
  citecolor=blue!45!black,urlcolor=blue!45!black]{hyperref}
\hypersetup{
  pdftitle={Exact W1 Central-Limit Asymptotics},
  pdfauthor={Gabriel Peyre}
}

\newtheorem{theorem}{Theorem}[section]
\newtheorem{proposition}[theorem]{Proposition}
\newtheorem{lemma}[theorem]{Lemma}
\theoremstyle{remark}
\newtheorem{remark}[theorem]{Remark}

\newcommand{\R}{\mathbb{R}}
\newcommand{\E}{\mathbb{E}}
\newcommand{\Wass}{\mathcal{W}}
\newcommand{\dd}{\,\mathrm{d}}
\newcommand{\abs}[1]{\left\lvert #1\right\rvert}
\newcommand{\Law}{\operatorname{Law}}
\newcommand{\Unif}{\operatorname{Unif}}

\title{Exact $\mathcal W_1$ Central-Limit Asymptotics\\
  \large Bernoulli and continuous uniform input laws}
\author{Gabriel Peyr\'e}
\date{July 2026}

\begin{document}
\maketitle

\begin{abstract}
The Berry--Esseen figure in OT4ML compares the standard Gaussian with
normalized sums generated by two centered, unit-variance laws: the symmetric
Bernoulli law and the continuous uniform law on a centered interval. The
second law is interpreted theoretically as $\Unif[-\sqrt3,\sqrt3]$; the finite
grid used in the notebook is only a numerical quadrature of this continuous
distribution. The two laws have different sharp asymptotics. The Bernoulli
law retains a lattice continuity correction of order $n^{-1/2}$, whereas
symmetry cancels the first smooth Edgeworth term for the non-lattice uniform
law, whose leading error is of order $n^{-1}$. More precisely,
\[
  \Wass_1(\alpha_n,\mathcal N(0,1))
  =\frac{1}{2\sqrt n}+o(n^{-1/2})
\]
for Bernoulli input, while for continuous uniform input,
\[
  \Wass_1(\alpha_n,\mathcal N(0,1))
  =\frac{1+4e^{-3/2}}{10\sqrt{2\pi}}\frac1n+o(n^{-1})
  =\frac{0.0755006500\ldots}{n}+o(n^{-1}).
\]
\end{abstract}

\section{Setup and exact answers}

Let $X_1,X_2,\ldots$ be independent copies of a centered random variable
$X$ with unit variance, and set
\begin{equation}\label{eq:normalized-sum}
  Z_n\coloneqq \frac{1}{\sqrt n}\sum_{i=1}^n X_i,
  \qquad
  \alpha_n\coloneqq\Law(Z_n),
  \qquad
  \gamma\coloneqq\mathcal N(0,1).
\end{equation}
This agrees with the numerical experiment, where
$\alpha_n=(D_{1/\sqrt n})_\sharp\alpha_0^{*n}$. We denote the standard normal
density and CDF by
\begin{equation}\label{eq:normal-density-cdf}
  \varphi(x)=\frac{e^{-x^2/2}}{\sqrt{2\pi}},
  \qquad
  \Phi(x)=\int_{-\infty}^x\varphi(t)\dd t.
\end{equation}

The calculations below give the following exact leading terms.
\begin{center}
\begin{tabular}{@{}lll@{}}
\toprule
Input law & Mechanism & Exact $\mathcal W_1$ asymptotic \\
\midrule
$\frac12(\delta_{-1}+\delta_{+1})$
  & lattice span $h=2$
  & $\frac{1}{2\sqrt n}+o(n^{-1/2})$ \\[1mm]
$\Unif[-\sqrt3,\sqrt3]$
  & fourth cumulant $\kappa_4=-6/5$
  & $\frac{1+4e^{-3/2}}{10\sqrt{2\pi}}\frac1n+o(n^{-1})$ \\
\bottomrule
\end{tabular}
\end{center}

The exponent changes because the first law is lattice-valued and the second
is genuinely continuous. Symmetry alone does not improve the Bernoulli rate:
it removes the smooth skewness term but not the oscillatory continuity
correction created by the lattice.

\section{The Bernoulli law: a lattice correction}

\paragraph{Maximal span.}
A nondegenerate law is lattice-valued if its support is contained in
$a+h\mathbb Z$ for some $a\in\R$ and $h>0$. Its \emph{maximal span} is the
largest admissible $h$. If $X\in a+h\mathbb Z$, then $Z_n$ lies on a lattice
whose spacing is
\begin{equation}\label{eq:normalized-span}
  \Delta_n=\frac{h}{\sqrt n}.
\end{equation}
This shrinking spacing leaves a first-order continuity correction.

Define the centered, one-periodic sawtooth
\begin{equation}\label{eq:sawtooth}
  \psi(u)=\frac12-\{u\},
\end{equation}
where $\{u\}=u-\lfloor u\rfloor$. Its value at integers is immaterial for
Lebesgue integration, and
\begin{equation}\label{eq:sawtooth-average}
  \int_0^1\abs{\psi(u)}\dd u=\frac14.
\end{equation}

\begin{lemma}[Periodic averaging]\label{lem:periodic-averaging}
Let $q$ be a bounded, Riemann-integrable, one-periodic function, let
$f\in L^1(\R)$, and let $c_\lambda\in\R$ be arbitrary. Then
\begin{equation}\label{eq:periodic-averaging}
  \lim_{\lambda\to+\infty}
  \int_{\R}f(x)q(\lambda x+c_\lambda)\dd x
  =
  \left(\int_{\R}f(x)\dd x\right)
  \left(\int_0^1q(u)\dd u\right).
\end{equation}
\end{lemma}

\begin{proof}
First suppose that $f$ is continuous and compactly supported. Partition its
support into intervals of length $1/\lambda$. Uniform continuity permits
replacing $f$ on each interval by its value at one endpoint, while periodicity
makes the corresponding integral of $q(\lambda x+c_\lambda)$ equal to
$\lambda^{-1}\int_0^1q$, apart from at most two boundary intervals. The
result is a Riemann sum for $\int f$, uniformly in the phase $c_\lambda$.
Approximation in $L^1$ and boundedness of $q$ give the general case.
\end{proof}

The classical lattice Edgeworth expansion supplies the oscillatory term. We
state only the integrated form needed here. It follows from standard
nonuniform lattice expansions; finite-support laws satisfy all the required
moment and remainder assumptions. See Petrov~\cite[Chapter~VII]{Petrov1975},
Bhattacharya and Ranga Rao~\cite[Chapter~5]{BhattacharyaRao2010}, and the
continuity-corrected construction of Kolassa and
McCullagh~\cite{KolassaMcCullagh1990}.

\begin{proposition}[Integrated lattice Edgeworth expansion]
\label{prop:lattice-edgeworth}
Let $X$ be centered, symmetric, of unit variance, and supported on a finite
subset of $a+h\mathbb Z$, where $h$ is the maximal span. If $F_n$ denotes the
CDF of $Z_n$, then
\begin{equation}\label{eq:lattice-edgeworth}
  F_n(x)-\Phi(x)
  =
  \frac{h}{\sqrt n}\,
  \psi\!\left(\frac{\sqrt n\,x-na}{h}\right)\varphi(x)
  +r_n(x),
\end{equation}
where
\begin{equation}\label{eq:lattice-edgeworth-remainder}
  \int_{\R}\abs{r_n(x)}\dd x=o(n^{-1/2}).
\end{equation}
\end{proposition}

For probability measures on the line with finite first moments, Vallender's
identity~\cite{Vallender1974} reads
\begin{equation}\label{eq:w1-cdf}
  \Wass_1(\alpha_n,\gamma)
  =\int_{\R}\abs{F_n(x)-\Phi(x)}\dd x.
\end{equation}
The sharp lattice constant follows by averaging the absolute continuity
correction.

\begin{theorem}[Sharp symmetric-lattice constant]
\label{thm:sharp-lattice-w1}
Under the assumptions of Proposition~\ref{prop:lattice-edgeworth},
\begin{equation}\label{eq:sharp-lattice-w1}
  \lim_{n\to\infty}\sqrt n\,\Wass_1(\alpha_n,\gamma)=\frac h4.
\end{equation}
\end{theorem}

\begin{proof}
Combining \eqref{eq:lattice-edgeworth} and \eqref{eq:w1-cdf}, and using
$\bigl|\abs{u+v}-\abs{u}\bigr|\leq\abs{v}$, gives
\begin{align}
  &\left|
  \sqrt n\,\Wass_1(\alpha_n,\gamma)
  -h\int_{\R}
  \abs{\psi\!\left(\frac{\sqrt n\,x-na}{h}\right)}
  \varphi(x)\dd x
  \right| \notag\\
  &\hspace{7em}\leq
  \sqrt n\int_{\R}\abs{r_n(x)}\dd x
  \longrightarrow0.
  \label{eq:w1-remainder-control}
\end{align}
Lemma~\ref{lem:periodic-averaging}, applied with $f=\varphi$ and
$q=\abs{\psi}$, shows that the remaining integral converges to $1/4$.
\end{proof}

For the symmetric Bernoulli law, the support is contained in
$-1+2\mathbb Z$, and its maximal span is $h=2$. Consequently,
\begin{equation}\label{eq:bernoulli-final}
  \boxed{
  \Wass_1\!\left(
    \Law\!\left(\frac1{\sqrt n}\sum_{i=1}^nX_i\right),
    \mathcal N(0,1)
  \right)
  =\frac1{2\sqrt n}+o(n^{-1/2}).}
\end{equation}
For example, at $n=160$ the leading term is
$1/(2\sqrt{160})=0.03952847\ldots$.

\section{The continuous uniform law: a smooth correction}

The second theoretical input is uniform on a centered interval. If
$X\sim\Unif[-a,a]$, then $\operatorname{Var}(X)=a^2/3$; unit variance fixes
$a=\sqrt3$. This law is non-lattice, symmetric, compactly supported, and its
characteristic function satisfies Cram\'er's condition. Its first nonzero
Edgeworth correction is therefore the smooth fourth-cumulant term.

\begin{proposition}[Symmetric non-lattice $\mathcal W_1$ asymptotic]
\label{prop:nonlattice-w1}
Let $X$ be centered, symmetric, and of unit variance. Assume that
$\E[\abs{X}^m]<+\infty$ for every $m\geq1$ and that Cram\'er's condition holds:
\begin{equation}\label{eq:cramer-condition}
  \limsup_{\abs{t}\to+\infty}\abs{\E[e^{itX}]}<1.
\end{equation}
With fourth cumulant
$\kappa_4=\E[X^4]-3$, one has
\begin{equation}\label{eq:nonlattice-edgeworth}
  F_n(x)-\Phi(x)
  =-\frac{\kappa_4}{24n}H_3(x)\varphi(x)+r_n(x),
  \qquad
  H_3(x)=x^3-3x,
\end{equation}
where $\int_\R\abs{r_n(x)}\dd x=o(n^{-1})$. Consequently,
\begin{equation}\label{eq:nonlattice-w1}
  \Wass_1(\alpha_n,\gamma)
  =\frac{\abs{\kappa_4}}{24n}
  \int_{\R}\abs{H_3(x)}\varphi(x)\dd x
  +o(n^{-1}).
\end{equation}
\end{proposition}

\begin{proof}
The non-lattice Edgeworth expansion through order $n^{-1}$ contains a term
of order $n^{-1/2}$ proportional to the third cumulant and terms of order
$n^{-1}$ proportional to the fourth cumulant and to the square of the third
cumulant. Symmetry makes the third cumulant vanish, leaving
\eqref{eq:nonlattice-edgeworth}; see
Bhattacharya and Ranga Rao~\cite[Chapter~4]{BhattacharyaRao2010} or
Petrov~\cite[Chapter~VI]{Petrov1975}. Inserting this expansion into
Vallender's identity \eqref{eq:w1-cdf} and using
$\bigl|\abs{u+v}-\abs u\bigr|\leq\abs v$ gives
\eqref{eq:nonlattice-w1}.
\end{proof}

It remains to evaluate the Hermite integral. Since
$H_3(x)=x(x^2-3)$ changes sign at $0$ and $\pm\sqrt3$, while
\begin{equation}\label{eq:hermite-antiderivative}
  \frac{\dd}{\dd x}\left[-(x^2-1)\varphi(x)\right]
  =H_3(x)\varphi(x),
\end{equation}
symmetry and direct integration give
\begin{equation}\label{eq:h3-integral}
  I_3\coloneqq\int_{\R}\abs{H_3(x)}\varphi(x)\dd x
  =2\varphi(0)+8\varphi(\sqrt3)
  =\frac{2+8e^{-3/2}}{\sqrt{2\pi}}.
\end{equation}

For $X\sim\Unif[-\sqrt3,\sqrt3]$,
\begin{equation}\label{eq:uniform-cumulant}
  \E[X^4]=\frac95,
  \qquad
  \kappa_4=\E[X^4]-3=-\frac65.
\end{equation}
Proposition~\ref{prop:nonlattice-w1} and \eqref{eq:h3-integral} therefore
yield the exact asymptotic
\begin{equation}\label{eq:uniform-final}
  \boxed{
  \Wass_1\!\left(
    \Law\!\left(\frac1{\sqrt n}\sum_{i=1}^nX_i\right),
    \mathcal N(0,1)
  \right)
  =\frac{1+4e^{-3/2}}{10\sqrt{2\pi}}\frac1n+o(n^{-1}).}
\end{equation}
Numerically,
\begin{equation}\label{eq:uniform-numerical-constant}
  \frac{1+4e^{-3/2}}{10\sqrt{2\pi}}
  =0.0755006500\ldots.
\end{equation}
At $n=160$, the corresponding leading prediction is
$0.0004718791\ldots$.

\section{Interpretation for the numerical figure}

The two theoretical rates are genuinely different:
\begin{equation}\label{eq:comparison-final}
  \Wass_1(\alpha_n,\gamma)
  \sim
  \begin{cases}
    \dfrac{1}{2\sqrt n},
      & \alpha_0=\frac12(\delta_{-1}+\delta_{+1}),\\[2mm]
    \dfrac{1+4e^{-3/2}}{10\sqrt{2\pi}}\dfrac1n,
      & \alpha_0=\Unif[-\sqrt3,\sqrt3].
  \end{cases}
\end{equation}
Thus the continuous uniform example displays the improved $n^{-1}$ rate
caused by the simultaneous absence of lattice oscillations and cancellation
of the skewness term. The 101-point grid in the notebook should be understood
as a quadrature rule for this continuous law, not as the law whose asymptotic
$n\to\infty$ behavior is being studied. Its resolution is therefore a
numerical parameter that should be refined until the discretization error is
negligible relative to the continuous $n^{-1}$ signal.

\begin{thebibliography}{9}

\bibitem{BhattacharyaRao2010}
R.~N. Bhattacharya and R.~Ranga Rao.
\newblock \emph{Normal Approximation and Asymptotic Expansions}.
\newblock Classics in Applied Mathematics 64, SIAM, Philadelphia, 2010.

\bibitem{KolassaMcCullagh1990}
J.~E. Kolassa and P.~McCullagh.
\newblock Edgeworth series for lattice distributions.
\newblock \emph{The Annals of Statistics}, 18(2):981--985, 1990.
\newblock \url{https://galton.uchicago.edu/~pmcc/pubs/paper16.pdf}.

\bibitem{Petrov1975}
V.~V. Petrov.
\newblock \emph{Sums of Independent Random Variables}.
\newblock Ergebnisse der Mathematik und ihrer Grenzgebiete 82,
Springer-Verlag, Berlin, 1975.

\bibitem{Vallender1974}
S.~S. Vallender.
\newblock Calculation of the Wasserstein distance between probability
distributions on the line.
\newblock \emph{Theory of Probability and Its Applications}, 18(4):784--786,
1974.
\newblock \url{https://doi.org/10.1137/1118101}.

\end{thebibliography}

\end{document}
